% !TEX root = ../main.tex

\textit{Mech Lab Tools} è una toolbox in Python che sto sviluppando per le esperienze di laboratorio del primo anno di fisica. L'obiettivo è fornire una base modulare e riutilizzabile per le operazioni ricorrenti nell'analisi dei dati sperimentali: statistica descrittiva e produzione di grafici.

Il progetto è organizzato in moduli distinti, ciascuno con un compito preciso, in modo che le stesse funzioni possano essere impiegate in esperienze diverse senza duplicazione di codice. I criteri guida dello sviluppo sono la trasparenza dei passaggi matematici, la riproducibilità dei risultati e la crescita graduale della toolbox man mano che le esigenze del corso aumentano.

Questa documentazione descrive ogni modulo in dettaglio — le funzioni implementate, le formule matematiche sottostanti, i parametri accettati e i valori restituiti — e costituisce un riferimento sia per l'utilizzo corretto degli strumenti sia per la verifica della correttezza delle implementazioni.


Le librerie esterne utilizzate nel progetto sono \texttt{NumPy}, \texttt{pandas} e \texttt{Matplotlib}. In tutto il codice vengono importate con i seguenti acronimi standard:
\begin{minted}{python}
import numpy as np
import pandas as pd
import matplotlib.pyplot as plt
\end{minted}

Gli snippet di codice sono formattati usando lo stile \texttt{gruvbox-light}.

Una riga ricorrente all'inizio di quasi tutte le funzioni è:
\begin{minted}{python}
x = np.asarray(x, dtype=float)
\end{minted}
Questa istruzione converte l'argomento \texttt{x} in un array NumPy di numeri in virgola mobile (\texttt{float64}), indipendentemente dal tipo con cui viene passato (lista, tupla, array intero, ecc.). Ciò garantisce che le operazioni vettoriali di NumPy funzionino correttamente e previene errori di tipo nelle divisioni o nelle operazioni matematiche. È una pratica difensiva comune nelle librerie numeriche in Python.
\hfill
